% Documentation of AI Usage (Appendix)
%%%%%%%%%%%%%%%%%%%%%%%%%%%%%%%%%%%%%%%%%%%%%%%%%%%%%%%%%
% Based on the LUH guideline "KI-Tools richtig dokumentieren"
% https://www.uni-hannover.de/fileadmin/luh/content/online-lehre/LUH_Handreichung_Dokumentation_KI-Einsatz.pdf
%
% IMPORTANT: Your supervisor decides how AI usage must be documented
% (e.g. whether prompts or chat logs have to be included). Check beforehand!
%
% If you did NOT use any AI tools:
%   - comment out % Documentation of AI Usage (Appendix)
%%%%%%%%%%%%%%%%%%%%%%%%%%%%%%%%%%%%%%%%%%%%%%%%%%%%%%%%%
% Based on the LUH guideline "KI-Tools richtig dokumentieren"
% https://www.uni-hannover.de/fileadmin/luh/content/online-lehre/LUH_Handreichung_Dokumentation_KI-Einsatz.pdf
%
% IMPORTANT: Your supervisor decides how AI usage must be documented
% (e.g. whether prompts or chat logs have to be included). Check beforehand!
%
% If you did NOT use any AI tools:
%   - comment out % Documentation of AI Usage (Appendix)
%%%%%%%%%%%%%%%%%%%%%%%%%%%%%%%%%%%%%%%%%%%%%%%%%%%%%%%%%
% Based on the LUH guideline "KI-Tools richtig dokumentieren"
% https://www.uni-hannover.de/fileadmin/luh/content/online-lehre/LUH_Handreichung_Dokumentation_KI-Einsatz.pdf
%
% IMPORTANT: Your supervisor decides how AI usage must be documented
% (e.g. whether prompts or chat logs have to be included). Check beforehand!
%
% If you did NOT use any AI tools:
%   - comment out % Documentation of AI Usage (Appendix)
%%%%%%%%%%%%%%%%%%%%%%%%%%%%%%%%%%%%%%%%%%%%%%%%%%%%%%%%%
% Based on the LUH guideline "KI-Tools richtig dokumentieren"
% https://www.uni-hannover.de/fileadmin/luh/content/online-lehre/LUH_Handreichung_Dokumentation_KI-Einsatz.pdf
%
% IMPORTANT: Your supervisor decides how AI usage must be documented
% (e.g. whether prompts or chat logs have to be included). Check beforehand!
%
% If you did NOT use any AI tools:
%   - comment out \input{91_ai_documentation.tex} in 00_main.tex
%   - remove the reference to this appendix in 03_declaration.tex
%
% Table columns (following the LUH guideline):
%   No.              consecutive number (to reference from footnotes/chat logs)
%   AI tool          tool name and provider, incl. link (model/version if known)
%   Purpose          short description of what the tool was used for
%   Affected parts   which sections or whole chapters were affected
%   Date & prompt    date of use; prompt depending on purpose/supervisor's requirements
%   Remarks          depth of processing or pointer to more detailed documentation
%
% The rows below are examples — delete tools you did not use, add your own.
%
% Further sources (cited in the guideline):
%   - Universität Basel, Vizerektorat Lehre: "Leitfaden 'aus KI' zitieren", June 2024
%   - Baresel et al.: "KI-Gebrauch im Studienkontext dokumentieren", 29 Nov 2024
%   - Frankfurt UAS: "Lehren, Lernen und Schreiben mit generativen KI-Anwendungen", 11 Mar 2024
%%%%%%%%%%%%%%%%%%%%%%%%%%%%%%%%%%%%%%%%%%%%%%%%%%%%%%%%%

\chapter{Documentation of AI Usage}
\label{cha:AI-Documentation}

AI tools were used as aids in this thesis. To keep the origin of wording and content traceable, their use is documented in table~\ref{tab:AI-Usage}. All AI-generated content was reviewed by me; I bear full responsibility for the content of this thesis.
% Optional, if chat logs are documented (uncomment the section below):
% Selected chat logs are reproduced in section~\ref{sec:AI-ChatLogs}.

% Optional: additionally mark AI usage in the text with a footnote, e.g.:
%   ... text.\footnote{Language edited with ChatGPT (OpenAI), 12.03.2026; see table~\ref{tab:AI-Usage}, no.~3.}

{%
\footnotesize
\setstretch{1.15}
\setlength{\tabcolsep}{4pt}
\urlstyle{same}
\begin{longtable}{@{}
    L{\dimexpr0.05\linewidth-2\tabcolsep}
    L{\dimexpr0.22\linewidth-2\tabcolsep}
    L{\dimexpr0.18\linewidth-2\tabcolsep}
    L{\dimexpr0.15\linewidth-2\tabcolsep}
    L{\dimexpr0.22\linewidth-2\tabcolsep}
    L{\dimexpr0.18\linewidth-2\tabcolsep}@{}}
    \caption[Documentation of AI usage.]{Documentation of the AI tools used in this thesis (following the guideline of Leibniz University Hannover).}
    \label{tab:AI-Usage}\\
    \toprule
    \textbf{No.} & \textbf{AI tool} & \textbf{Purpose} & \textbf{Affected part(s) of the thesis} & \textbf{Date of use \& prompt} & \textbf{Remarks} \\
    \midrule
    \endfirsthead
    \caption[]{(continued) Documentation of the AI tools used in this thesis.}\\
    \toprule
    \textbf{No.} & \textbf{AI tool} & \textbf{Purpose} & \textbf{Affected part(s) of the thesis} & \textbf{Date of use \& prompt} & \textbf{Remarks} \\
    \midrule
    \endhead
    \midrule
    \multicolumn{6}{r@{}}{\textit{Continued on next page}} \\
    \endfoot
    \bottomrule
    \endlastfoot
    %--- Example rows (replace) -----------------------------------------
    1 & LUHKI (LUH)\newline \href{https://luhki.uni-hannover.de}{luhki.uni-hannover.de}
      & Creating an outline
      & entire thesis
      & 04.03.2026
      & Only ideas adopted, heavily modified \\
    \midrule
    2 & DeepL (institute licence)\newline \href{https://www.deepl.com}{deepl.com}
      & Translation of text passages
      & Chapter 2
      & 10.03.2026
      & Translation checked and revised \\
    \midrule
    3 & ChatGPT (OpenAI)\newline \href{https://chatgpt.com}{chatgpt.com}
      & Linguistic and formal revision
      & Chapter 3, section 3.2
      & 12.03.2026\newline \enquote{Check the spelling and punctuation in the following section}
      & Only corrections adopted \\
    \midrule
    4 & Claude Code (Anthropic)\newline \href{https://claude.com/claude-code}{claude.com/claude-code}
      & Programming of evaluation scripts (Python)
      & Chapter 4, appendix A
      & 18.03.2026
      & Code reviewed, tested and adapted \\
    \midrule
    5 & Gemini (Google)\newline \href{https://gemini.google.com}{gemini.google.com}
      & Literature search
      & Chapter 2, section 2.1
      & 20.03.2026
      & Initial overview; sources then checked myself \\
\end{longtable}%
}

%=== Chat logs (optional) ===============================
% Disabled by default. To enable, uncomment the block below
% (select all lines and remove the comment signs, in Overleaf: Ctrl+/).
% Create one subsection per chat log and state its number from the table.
% Put prompts in the 'aiprompt' box and responses in the 'airesponse' box.
% Long responses may be shortened — mark omissions with [...].
%
% Alternatively, include an exported chat log as PDF:
%   \includepdf[pages=-, scale=0.85, pagecommand={}]{figures/chatlog_no3.pdf}
%
% \section{Chat Logs}
% \label{sec:AI-ChatLogs}
%
% The following chat logs supplement the information in table~\ref{tab:AI-Usage}. Longer responses have been shortened; omissions are marked with [\,\dots].
%
% \subsection{No. 3: Revision of Section 3.2}
% \label{sec:AI-Chat-3}
%
% \textbf{Tool:} ChatGPT (OpenAI) \quad
% \textbf{Date:} 12.03.2026
%
% \begin{aiprompt}[Prompt 1]
% Rephrase the following paragraph in a more scientific and precise way without changing its content:
%
% \enquote{The rotor blades are heavily loaded at high wind speeds because the forces get bigger and the blades then bend more.}
% \end{aiprompt}
%
% \begin{airesponse}[Response 1]
% \enquote{With increasing wind speed, the aerodynamic loads on the rotor blades rise, resulting in larger blade deflections and thus increased structural stress.}
% \end{airesponse}
%
% \textit{Adoption:} Wording from response~1 adopted in essence and adapted.

\clearpage
 in 00_main.tex
%   - remove the reference to this appendix in 03_declaration.tex
%
% Table columns (following the LUH guideline):
%   No.              consecutive number (to reference from footnotes/chat logs)
%   AI tool          tool name and provider, incl. link (model/version if known)
%   Purpose          short description of what the tool was used for
%   Affected parts   which sections or whole chapters were affected
%   Date & prompt    date of use; prompt depending on purpose/supervisor's requirements
%   Remarks          depth of processing or pointer to more detailed documentation
%
% The rows below are examples — delete tools you did not use, add your own.
%
% Further sources (cited in the guideline):
%   - Universität Basel, Vizerektorat Lehre: "Leitfaden 'aus KI' zitieren", June 2024
%   - Baresel et al.: "KI-Gebrauch im Studienkontext dokumentieren", 29 Nov 2024
%   - Frankfurt UAS: "Lehren, Lernen und Schreiben mit generativen KI-Anwendungen", 11 Mar 2024
%%%%%%%%%%%%%%%%%%%%%%%%%%%%%%%%%%%%%%%%%%%%%%%%%%%%%%%%%

\chapter{Documentation of AI Usage}
\label{cha:AI-Documentation}

AI tools were used as aids in this thesis. To keep the origin of wording and content traceable, their use is documented in table~\ref{tab:AI-Usage}. All AI-generated content was reviewed by me; I bear full responsibility for the content of this thesis.
% Optional, if chat logs are documented (uncomment the section below):
% Selected chat logs are reproduced in section~\ref{sec:AI-ChatLogs}.

% Optional: additionally mark AI usage in the text with a footnote, e.g.:
%   ... text.\footnote{Language edited with ChatGPT (OpenAI), 12.03.2026; see table~\ref{tab:AI-Usage}, no.~3.}

{%
\footnotesize
\setstretch{1.15}
\setlength{\tabcolsep}{4pt}
\urlstyle{same}
\begin{longtable}{@{}
    L{\dimexpr0.05\linewidth-2\tabcolsep}
    L{\dimexpr0.22\linewidth-2\tabcolsep}
    L{\dimexpr0.18\linewidth-2\tabcolsep}
    L{\dimexpr0.15\linewidth-2\tabcolsep}
    L{\dimexpr0.22\linewidth-2\tabcolsep}
    L{\dimexpr0.18\linewidth-2\tabcolsep}@{}}
    \caption[Documentation of AI usage.]{Documentation of the AI tools used in this thesis (following the guideline of Leibniz University Hannover).}
    \label{tab:AI-Usage}\\
    \toprule
    \textbf{No.} & \textbf{AI tool} & \textbf{Purpose} & \textbf{Affected part(s) of the thesis} & \textbf{Date of use \& prompt} & \textbf{Remarks} \\
    \midrule
    \endfirsthead
    \caption[]{(continued) Documentation of the AI tools used in this thesis.}\\
    \toprule
    \textbf{No.} & \textbf{AI tool} & \textbf{Purpose} & \textbf{Affected part(s) of the thesis} & \textbf{Date of use \& prompt} & \textbf{Remarks} \\
    \midrule
    \endhead
    \midrule
    \multicolumn{6}{r@{}}{\textit{Continued on next page}} \\
    \endfoot
    \bottomrule
    \endlastfoot
    %--- Example rows (replace) -----------------------------------------
    1 & LUHKI (LUH)\newline \href{https://luhki.uni-hannover.de}{luhki.uni-hannover.de}
      & Creating an outline
      & entire thesis
      & 04.03.2026
      & Only ideas adopted, heavily modified \\
    \midrule
    2 & DeepL (institute licence)\newline \href{https://www.deepl.com}{deepl.com}
      & Translation of text passages
      & Chapter 2
      & 10.03.2026
      & Translation checked and revised \\
    \midrule
    3 & ChatGPT (OpenAI)\newline \href{https://chatgpt.com}{chatgpt.com}
      & Linguistic and formal revision
      & Chapter 3, section 3.2
      & 12.03.2026\newline \enquote{Check the spelling and punctuation in the following section}
      & Only corrections adopted \\
    \midrule
    4 & Claude Code (Anthropic)\newline \href{https://claude.com/claude-code}{claude.com/claude-code}
      & Programming of evaluation scripts (Python)
      & Chapter 4, appendix A
      & 18.03.2026
      & Code reviewed, tested and adapted \\
    \midrule
    5 & Gemini (Google)\newline \href{https://gemini.google.com}{gemini.google.com}
      & Literature search
      & Chapter 2, section 2.1
      & 20.03.2026
      & Initial overview; sources then checked myself \\
\end{longtable}%
}

%=== Chat logs (optional) ===============================
% Disabled by default. To enable, uncomment the block below
% (select all lines and remove the comment signs, in Overleaf: Ctrl+/).
% Create one subsection per chat log and state its number from the table.
% Put prompts in the 'aiprompt' box and responses in the 'airesponse' box.
% Long responses may be shortened — mark omissions with [...].
%
% Alternatively, include an exported chat log as PDF:
%   \includepdf[pages=-, scale=0.85, pagecommand={}]{figures/chatlog_no3.pdf}
%
% \section{Chat Logs}
% \label{sec:AI-ChatLogs}
%
% The following chat logs supplement the information in table~\ref{tab:AI-Usage}. Longer responses have been shortened; omissions are marked with [\,\dots].
%
% \subsection{No. 3: Revision of Section 3.2}
% \label{sec:AI-Chat-3}
%
% \textbf{Tool:} ChatGPT (OpenAI) \quad
% \textbf{Date:} 12.03.2026
%
% \begin{aiprompt}[Prompt 1]
% Rephrase the following paragraph in a more scientific and precise way without changing its content:
%
% \enquote{The rotor blades are heavily loaded at high wind speeds because the forces get bigger and the blades then bend more.}
% \end{aiprompt}
%
% \begin{airesponse}[Response 1]
% \enquote{With increasing wind speed, the aerodynamic loads on the rotor blades rise, resulting in larger blade deflections and thus increased structural stress.}
% \end{airesponse}
%
% \textit{Adoption:} Wording from response~1 adopted in essence and adapted.

\clearpage
 in 00_main.tex
%   - remove the reference to this appendix in 03_declaration.tex
%
% Table columns (following the LUH guideline):
%   No.              consecutive number (to reference from footnotes/chat logs)
%   AI tool          tool name and provider, incl. link (model/version if known)
%   Purpose          short description of what the tool was used for
%   Affected parts   which sections or whole chapters were affected
%   Date & prompt    date of use; prompt depending on purpose/supervisor's requirements
%   Remarks          depth of processing or pointer to more detailed documentation
%
% The rows below are examples — delete tools you did not use, add your own.
%
% Further sources (cited in the guideline):
%   - Universität Basel, Vizerektorat Lehre: "Leitfaden 'aus KI' zitieren", June 2024
%   - Baresel et al.: "KI-Gebrauch im Studienkontext dokumentieren", 29 Nov 2024
%   - Frankfurt UAS: "Lehren, Lernen und Schreiben mit generativen KI-Anwendungen", 11 Mar 2024
%%%%%%%%%%%%%%%%%%%%%%%%%%%%%%%%%%%%%%%%%%%%%%%%%%%%%%%%%

\chapter{Documentation of AI Usage}
\label{cha:AI-Documentation}

AI tools were used as aids in this thesis. To keep the origin of wording and content traceable, their use is documented in table~\ref{tab:AI-Usage}. All AI-generated content was reviewed by me; I bear full responsibility for the content of this thesis.
% Optional, if chat logs are documented (uncomment the section below):
% Selected chat logs are reproduced in section~\ref{sec:AI-ChatLogs}.

% Optional: additionally mark AI usage in the text with a footnote, e.g.:
%   ... text.\footnote{Language edited with ChatGPT (OpenAI), 12.03.2026; see table~\ref{tab:AI-Usage}, no.~3.}

{%
\footnotesize
\setstretch{1.15}
\setlength{\tabcolsep}{4pt}
\urlstyle{same}
\begin{longtable}{@{}
    L{\dimexpr0.05\linewidth-2\tabcolsep}
    L{\dimexpr0.22\linewidth-2\tabcolsep}
    L{\dimexpr0.18\linewidth-2\tabcolsep}
    L{\dimexpr0.15\linewidth-2\tabcolsep}
    L{\dimexpr0.22\linewidth-2\tabcolsep}
    L{\dimexpr0.18\linewidth-2\tabcolsep}@{}}
    \caption[Documentation of AI usage.]{Documentation of the AI tools used in this thesis (following the guideline of Leibniz University Hannover).}
    \label{tab:AI-Usage}\\
    \toprule
    \textbf{No.} & \textbf{AI tool} & \textbf{Purpose} & \textbf{Affected part(s) of the thesis} & \textbf{Date of use \& prompt} & \textbf{Remarks} \\
    \midrule
    \endfirsthead
    \caption[]{(continued) Documentation of the AI tools used in this thesis.}\\
    \toprule
    \textbf{No.} & \textbf{AI tool} & \textbf{Purpose} & \textbf{Affected part(s) of the thesis} & \textbf{Date of use \& prompt} & \textbf{Remarks} \\
    \midrule
    \endhead
    \midrule
    \multicolumn{6}{r@{}}{\textit{Continued on next page}} \\
    \endfoot
    \bottomrule
    \endlastfoot
    %--- Example rows (replace) -----------------------------------------
    1 & LUHKI (LUH)\newline \href{https://luhki.uni-hannover.de}{luhki.uni-hannover.de}
      & Creating an outline
      & entire thesis
      & 04.03.2026
      & Only ideas adopted, heavily modified \\
    \midrule
    2 & DeepL (institute licence)\newline \href{https://www.deepl.com}{deepl.com}
      & Translation of text passages
      & Chapter 2
      & 10.03.2026
      & Translation checked and revised \\
    \midrule
    3 & ChatGPT (OpenAI)\newline \href{https://chatgpt.com}{chatgpt.com}
      & Linguistic and formal revision
      & Chapter 3, section 3.2
      & 12.03.2026\newline \enquote{Check the spelling and punctuation in the following section}
      & Only corrections adopted \\
    \midrule
    4 & Claude Code (Anthropic)\newline \href{https://claude.com/claude-code}{claude.com/claude-code}
      & Programming of evaluation scripts (Python)
      & Chapter 4, appendix A
      & 18.03.2026
      & Code reviewed, tested and adapted \\
    \midrule
    5 & Gemini (Google)\newline \href{https://gemini.google.com}{gemini.google.com}
      & Literature search
      & Chapter 2, section 2.1
      & 20.03.2026
      & Initial overview; sources then checked myself \\
\end{longtable}%
}

%=== Chat logs (optional) ===============================
% Disabled by default. To enable, uncomment the block below
% (select all lines and remove the comment signs, in Overleaf: Ctrl+/).
% Create one subsection per chat log and state its number from the table.
% Put prompts in the 'aiprompt' box and responses in the 'airesponse' box.
% Long responses may be shortened — mark omissions with [...].
%
% Alternatively, include an exported chat log as PDF:
%   \includepdf[pages=-, scale=0.85, pagecommand={}]{figures/chatlog_no3.pdf}
%
% \section{Chat Logs}
% \label{sec:AI-ChatLogs}
%
% The following chat logs supplement the information in table~\ref{tab:AI-Usage}. Longer responses have been shortened; omissions are marked with [\,\dots].
%
% \subsection{No. 3: Revision of Section 3.2}
% \label{sec:AI-Chat-3}
%
% \textbf{Tool:} ChatGPT (OpenAI) \quad
% \textbf{Date:} 12.03.2026
%
% \begin{aiprompt}[Prompt 1]
% Rephrase the following paragraph in a more scientific and precise way without changing its content:
%
% \enquote{The rotor blades are heavily loaded at high wind speeds because the forces get bigger and the blades then bend more.}
% \end{aiprompt}
%
% \begin{airesponse}[Response 1]
% \enquote{With increasing wind speed, the aerodynamic loads on the rotor blades rise, resulting in larger blade deflections and thus increased structural stress.}
% \end{airesponse}
%
% \textit{Adoption:} Wording from response~1 adopted in essence and adapted.

\clearpage
 in 00_main.tex
%   - remove the reference to this appendix in 03_declaration.tex
%
% Table columns (following the LUH guideline):
%   No.              consecutive number (to reference from footnotes/chat logs)
%   AI tool          tool name and provider, incl. link (model/version if known)
%   Purpose          short description of what the tool was used for
%   Affected parts   which sections or whole chapters were affected
%   Date & prompt    date of use; prompt depending on purpose/supervisor's requirements
%   Remarks          depth of processing or pointer to more detailed documentation
%
% The rows below are examples — delete tools you did not use, add your own.
%
% Further sources (cited in the guideline):
%   - Universität Basel, Vizerektorat Lehre: "Leitfaden 'aus KI' zitieren", June 2024
%   - Baresel et al.: "KI-Gebrauch im Studienkontext dokumentieren", 29 Nov 2024
%   - Frankfurt UAS: "Lehren, Lernen und Schreiben mit generativen KI-Anwendungen", 11 Mar 2024
%%%%%%%%%%%%%%%%%%%%%%%%%%%%%%%%%%%%%%%%%%%%%%%%%%%%%%%%%

\chapter{Documentation of AI Usage}
\label{cha:AI-Documentation}

AI tools were used as aids in this thesis. To keep the origin of wording and content traceable, their use is documented in table~\ref{tab:AI-Usage}. All AI-generated content was reviewed by me; I bear full responsibility for the content of this thesis.
% Optional, if chat logs are documented (uncomment the section below):
% Selected chat logs are reproduced in section~\ref{sec:AI-ChatLogs}.

% Optional: additionally mark AI usage in the text with a footnote, e.g.:
%   ... text.\footnote{Language edited with ChatGPT (OpenAI), 12.03.2026; see table~\ref{tab:AI-Usage}, no.~3.}

{%
\footnotesize
\setstretch{1.15}
\setlength{\tabcolsep}{4pt}
\urlstyle{same}
\begin{longtable}{@{}
    L{\dimexpr0.05\linewidth-2\tabcolsep}
    L{\dimexpr0.22\linewidth-2\tabcolsep}
    L{\dimexpr0.18\linewidth-2\tabcolsep}
    L{\dimexpr0.15\linewidth-2\tabcolsep}
    L{\dimexpr0.22\linewidth-2\tabcolsep}
    L{\dimexpr0.18\linewidth-2\tabcolsep}@{}}
    \caption[Documentation of AI usage.]{Documentation of the AI tools used in this thesis (following the guideline of Leibniz University Hannover).}
    \label{tab:AI-Usage}\\
    \toprule
    \textbf{No.} & \textbf{AI tool} & \textbf{Purpose} & \textbf{Affected part(s) of the thesis} & \textbf{Date of use \& prompt} & \textbf{Remarks} \\
    \midrule
    \endfirsthead
    \caption[]{(continued) Documentation of the AI tools used in this thesis.}\\
    \toprule
    \textbf{No.} & \textbf{AI tool} & \textbf{Purpose} & \textbf{Affected part(s) of the thesis} & \textbf{Date of use \& prompt} & \textbf{Remarks} \\
    \midrule
    \endhead
    \midrule
    \multicolumn{6}{r@{}}{\textit{Continued on next page}} \\
    \endfoot
    \bottomrule
    \endlastfoot
    %--- Example rows (replace) -----------------------------------------
    1 & LUHKI (LUH)\newline \href{https://luhki.uni-hannover.de}{luhki.uni-hannover.de}
      & Creating an outline
      & entire thesis
      & 04.03.2026
      & Only ideas adopted, heavily modified \\
    \midrule
    2 & DeepL (institute licence)\newline \href{https://www.deepl.com}{deepl.com}
      & Translation of text passages
      & Chapter 2
      & 10.03.2026
      & Translation checked and revised \\
    \midrule
    3 & ChatGPT (OpenAI)\newline \href{https://chatgpt.com}{chatgpt.com}
      & Linguistic and formal revision
      & Chapter 3, section 3.2
      & 12.03.2026\newline \enquote{Check the spelling and punctuation in the following section}
      & Only corrections adopted \\
    \midrule
    4 & Claude Code (Anthropic)\newline \href{https://claude.com/claude-code}{claude.com/claude-code}
      & Programming of evaluation scripts (Python)
      & Chapter 4, appendix A
      & 18.03.2026
      & Code reviewed, tested and adapted \\
    \midrule
    5 & Gemini (Google)\newline \href{https://gemini.google.com}{gemini.google.com}
      & Literature search
      & Chapter 2, section 2.1
      & 20.03.2026
      & Initial overview; sources then checked myself \\
\end{longtable}%
}

%=== Chat logs (optional) ===============================
% Disabled by default. To enable, uncomment the block below
% (select all lines and remove the comment signs, in Overleaf: Ctrl+/).
% Create one subsection per chat log and state its number from the table.
% Put prompts in the 'aiprompt' box and responses in the 'airesponse' box.
% Long responses may be shortened — mark omissions with [...].
%
% Alternatively, include an exported chat log as PDF:
%   \includepdf[pages=-, scale=0.85, pagecommand={}]{figures/chatlog_no3.pdf}
%
% \section{Chat Logs}
% \label{sec:AI-ChatLogs}
%
% The following chat logs supplement the information in table~\ref{tab:AI-Usage}. Longer responses have been shortened; omissions are marked with [\,\dots].
%
% \subsection{No. 3: Revision of Section 3.2}
% \label{sec:AI-Chat-3}
%
% \textbf{Tool:} ChatGPT (OpenAI) \quad
% \textbf{Date:} 12.03.2026
%
% \begin{aiprompt}[Prompt 1]
% Rephrase the following paragraph in a more scientific and precise way without changing its content:
%
% \enquote{The rotor blades are heavily loaded at high wind speeds because the forces get bigger and the blades then bend more.}
% \end{aiprompt}
%
% \begin{airesponse}[Response 1]
% \enquote{With increasing wind speed, the aerodynamic loads on the rotor blades rise, resulting in larger blade deflections and thus increased structural stress.}
% \end{airesponse}
%
% \textit{Adoption:} Wording from response~1 adopted in essence and adapted.

\clearpage
